\chapter{Theoretical Background}
\label{chap:theory}

This chapter presents the theoretical foundations of the proposed framework.
We begin with the governing equations for transient heat conduction and the
boundary conditions used to model water quenching.
We then describe the Finite Element Method, which serves as the reference
solver throughout this work.
Next, we introduce Physics-Informed Neural Networks and the techniques used
to stabilise their training.
We follow this with an overview of Neural Architecture Search and the three
search strategies compared in this thesis.
The chapter closes with a description of the RBniCS Thermal Fin benchmark,
which provides a standard test case for the framework.

\section{Transient Heat Conduction}
\label{sec:heat_theory}

The temperature field $T(\mathbf{x},t)$ inside a solid domain $\Omega$ during
cooling is governed by the transient heat conduction equation
\cite{xiao2010,wang2019quenching}:

\begin{equation}
\rho c_p \frac{\partial T}{\partial t}
= \nabla \cdot (k_T \nabla T),
\qquad \mathbf{x}\in\Omega,\; t>0.
\label{eq:heat_full}
\end{equation}

% [HOCA #1]: "$\rho c_p$ - ?"
% DUZELTME: rho, c_p, k_T ve alpha birimleriyle birlikte asagida tanimlandi.
Here $\rho$ is the material density ($\mathrm{kg\,m^{-3}}$) and $c_p$ is the
specific heat capacity at constant pressure ($\mathrm{J\,kg^{-1}K^{-1}}$).
The product $\rho c_p$ is the volumetric heat capacity
($\mathrm{J\,m^{-3}K^{-1}}$) and $k_T$ is the thermal conductivity
($\mathrm{W\,m^{-1}K^{-1}}$).
The thermal diffusivity is $\alpha = k_T / (\rho c_p)$
($\mathrm{m^2\,s^{-1}}$).

\subsection{Initial and Boundary Conditions}

A uniform initial temperature is assumed:
\begin{equation}
T(\mathbf{x},0) = T_0, \qquad \mathbf{x}\in\Omega,
\label{eq:ic}
\end{equation}

% [HOCA #2]: "Define T_{0}"
% DUZELTME: T_0 asagida tanimlanmistir.
where $T_0 = 540\,^\circ\mathrm{C}$ is the initial uniform temperature of
the casting at the start of quenching.

Water quenching is modelled using a convective (Robin) boundary condition on
the entire exposed surface $\partial\Omega$:
\begin{equation}
-k_T \nabla T \cdot \hat{\mathbf{n}} = h(T - T_w),
\qquad \mathbf{x}\in\partial\Omega,
\label{eq:robin}
\end{equation}

where $h$ is the convective heat transfer coefficient
($\mathrm{W\,m^{-2}K^{-1}}$), $T_w = 20\,^\circ\mathrm{C}$ is the quench
bath temperature, and $\hat{\mathbf{n}}$ is the outward unit normal to the
surface.
The Robin condition captures the resistance to heat transfer at the
solid-liquid interface.
It avoids the need to solve the full conjugate heat transfer problem
\cite{wilson1966,xiao2010}.

\subsection{Dimensionless Formulation}

It is convenient to define the dimensionless temperature
\begin{equation}
\theta(\mathbf{x},t) = \frac{T(\mathbf{x},t) - T_w}{T_0 - T_w},
\label{eq:theta}
\end{equation}

so that $\theta=1$ at the initial hot state and $\theta=0$ at the water
temperature.
The governing equation becomes
\begin{equation}
\frac{\partial\theta}{\partial t} = \alpha\nabla^2\theta,
\label{eq:heat_dim}
\end{equation}

with initial condition $\theta(\mathbf{x},0)=1$ and Robin boundary condition
$-\nabla\theta\cdot\hat{\mathbf{n}} = (h/k_T)\,\theta$.

\subsection{Biot and Fourier Numbers}

The Biot number $\mathrm{Bi} = hL/k_T$ characterises the ratio of convective
to conductive resistance, where $L$ is a characteristic length.
For the A356 material and geometries used here, $\mathrm{Bi}\gg1$, which means
internal conduction is the limiting factor rather than surface convection.
The Fourier number $\mathrm{Fo} = \alpha t/L^2$ measures the dimensionless
time elapsed.

\section{Finite Element Method for Heat Transfer}
\label{sec:fem_theory}

The Finite Element Method discretises the weak form of
Eq.~\eqref{eq:heat_full} over a mesh of finite elements
\cite{wilson1966,mackerle2003}.
For a linear element formulation with implicit time integration, each time step
requires solving the linear system:

\begin{equation}
\left(\mathbf{C} + \Delta t\,\mathbf{K}\right)\mathbf{T}^{n+1}
= \mathbf{C}\,\mathbf{T}^{n} + \Delta t\,\mathbf{f}^{n+1},
\label{eq:fem_system}
\end{equation}

where $\mathbf{C}$ is the capacity matrix, $\mathbf{K}$ is the conductance
matrix, $\mathbf{f}$ contains boundary contributions, and $\mathbf{T}^n$ is
the nodal temperature vector at time $t_n$.
For large meshes, assembling and solving this system at each time step is
the dominant cost.

% [HOCA #3]: "Did the paper also used the same method?"
% "the FEM reference solution is computed by a Crank-Nicolson finite-difference (FD)"
% DUZELTME: Mortensen et al. tam FEM kullanmistir; bu tezde yapilandirilmis
% grid uzerinde Crank-Nicolson FD referansi kullanildigi aciklandi.
In this thesis, we computed the reference temperature field using a
Crank--Nicolson finite-difference (FD) solver on a structured grid.
This is not the same as the full industrial FEM used by Mortensen
et al.~\cite{mortensen2026}.
Mortensen et al.\ use a commercial FEM solver on an unstructured mesh with
coupled thermomechanical equations.
The structured-grid FD solver reproduces only the thermal sub-problem on the
benchmark geometries.
It serves as the ground-truth reference for training and evaluating the
NAS-PINN.
Its computational cost is substantially lower than the full industrial FEM.
The internal time step is $\Delta t_{\mathrm{int}}=0.5\,\mathrm{s}$.
The Crank--Nicolson scheme is unconditionally stable.
We exported reference temperature fields every $\Delta t_a=1.5\,\mathrm{s}$.
This gives 21 snapshots from $t=0$ to $t=30\,\mathrm{s}$.
In a real deployment, each export corresponds to one expensive FEM solver call.
Reducing the number of required exports is the core objective of the MSWP
framework.

The complete solver workflow (domain discretisation, element assembly,
global system solution, and comparison against the NAS-PINN prediction)
is illustrated in Figure~\ref{fig:fem_solver_method} in
Chapter~\ref{chap:method}.

\section{Physics-Informed Neural Networks}
\label{sec:pinn_theory}

\subsection{Formulation}

A PINN approximates the solution of a partial differential equation (PDE)
by training a neural network $N_\phi(\mathbf{x},t)$ to minimise a
physics-based loss \cite{raissi2019}:

\begin{equation}
\mathcal{L} = w_\mathrm{pde}\mathcal{L}_\mathrm{pde}
            + w_\mathrm{ic}\mathcal{L}_\mathrm{ic}
            + w_\mathrm{bc}\mathcal{L}_\mathrm{bc},
\label{eq:pinn_loss}
\end{equation}

% [HOCA #4]: "Define all variables in this equation."
% DUZELTME: w_pde, w_ic, w_bc ve kayip terimleri asagida tanimlandi.
where $w_\mathrm{pde}$, $w_\mathrm{ic}$, and $w_\mathrm{bc}$ are positive
scalar penalty weights that control the relative importance of each term.
$\mathcal{L}_\mathrm{pde}$ is the PDE residual loss that enforces the
governing physics.
$\mathcal{L}_\mathrm{ic}$ is the initial condition loss.
$\mathcal{L}_\mathrm{bc}$ is the boundary condition loss.
Each term is a mean-squared residual evaluated at collocation points sampled
inside the domain, at the initial time, and on the boundary, respectively.

% [HOCA #5]: "How this term relates to the previous equations?"
% DUZELTME: N_phi, isi denklemindeki theta'nin yerini alir (Eq. heat_dim).
The PDE loss enforces the heat equation (Eq.~\eqref{eq:heat_dim}) on the
network output.
The network $N_\phi(\mathbf{x},t)$ plays the role of the dimensionless
temperature $\theta$.
If $N_\phi$ were the exact solution, the residual would be zero.
Formally:
\begin{equation}
\mathcal{L}_\mathrm{pde}
= \mathrm{MSE}\!\left(\frac{\partial N_\phi}{\partial t}
  - \alpha\nabla^2 N_\phi\right).
\label{eq:lpde}
\end{equation}

Here $\alpha$ is the thermal diffusivity from Eq.~\eqref{eq:heat_dim}.
We compute the derivatives via automatic differentiation.

By including physical laws directly in the loss, PINNs require far fewer
labelled data points than purely supervised models
\cite{karniadakis2021,cuomo2022}.

\subsection{Spectral Bias and Fourier Feature Embeddings}

Standard neural networks exhibit a \emph{spectral bias} \cite{rahaman2019}.
They learn low-frequency components first.
Fine spatial or temporal structures are learned slowly or not at all.
In heat transfer problems, steep temperature gradients near boundaries and
corners contain high-frequency content.
A standard PINN may underfit these regions.

Fourier feature embeddings \cite{tancik2020} address this by replacing the
raw spatial coordinates with a random Fourier mapping:
\begin{equation}
\gamma(\mathbf{x}) = \bigl[\cos(2\pi \mathbf{B}\mathbf{x}),\;
                            \sin(2\pi \mathbf{B}\mathbf{x})\bigr],
\end{equation}

% [HOCA #6]: "What are the dimensions of x and B?"
% DUZELTME: x in R^d, B in R^{m x d} olarak asagida tanimlandi.
where $\mathbf{x} \in \mathbb{R}^d$ is the spatial coordinate vector ($d=2$
or $3$ depending on the domain) and $\mathbf{B} \in \mathbb{R}^{m\times d}$
is a fixed random frequency matrix with $m$ rows, one per Fourier feature pair.
Entries of $\mathbf{B}$ are sampled from $\mathcal{N}(0,\sigma^2)$.
The scale $\sigma$ controls the frequency range of the features.
In this framework, Fourier features are applied for domains with steep
gradients.

\subsection{Self-Adaptive Loss Weighting}

The performance of a PINN depends critically on the balance between loss terms.
Poorly set weights cause some physics constraints to dominate.
Other constraints are then barely enforced.
McClenny and Braga-Neto \cite{mcclenny2023} proposed \emph{self-adaptive}
weights: learnable parameters $w_i$ that are updated during training to balance
the gradient magnitudes of each loss term.

In this framework, the boundary and endpoint loss weights are learnable, while
the PDE and IC weights are held fixed.
% [HOCA #7]: "Is endpoint same with boundary? Why? Justification needed."
% DUZELTME: "Endpoint" zamansal pencere sonu (t_end) oldugu, uzaysal sinir
% olmadigi aciklandi.
Here \emph{endpoint} refers to the temporal end of each prediction window
($t = t_\mathrm{end}$), not the spatial boundary.
An additional loss term penalises the deviation of the network output from the
FEM-provided temperature field at the final time point of the window.
This prevents the PINN from drifting away from the physical solution over the
window duration.

\subsection{Training for Long Time Horizons}

Several studies have noted that PINNs struggle to propagate information
forward in time over long intervals \cite{wang2021,krishnapriyan2021}.
The physics residual and the initial condition loss pull the network in
opposing directions across a long time span.
Satisfying the governing equation at late times conflicts with memory of the
initial state.
% [HOCA #8]: "What is gradient pathology? First time I'm seeing this term."
% DUZELTME: "Gradient pathology" tanimlandi.
This leads to \emph{gradient pathology} \cite{wang2022ntk,mcclenny2023}.
Different loss terms produce gradients of vastly different magnitudes during
back-propagation.
The optimiser effectively ignores the smaller-gradient terms.
Some physics constraints are then under-enforced.

The \emph{time decomposition} approach addresses this directly.
The full simulation interval is divided into shorter sub-problems.
Each sub-problem is solved independently with a fresh model initialised from
the known temperature field at the sub-problem boundary.
% [HOCA #9]: "What is Delta t_k? What is k and 1.5s? How did you find 1.5s?"
% DUZELTME: Delta t_k = k * Delta t_a aciklandi. 1.5s'nin 30 saniyelik
% quench suresini 20 esit araliga bolmekten geldigi belirtildi.
Each window covers a duration of $\Delta t_k = k \cdot \Delta t_a$, where
$k \in \{1,\ldots,5\}$ is the skip factor (Chapter~\ref{chap:intro}) and
$\Delta t_a = 1.5\,\mathrm{s}$ is the FEM anchor spacing.
The anchor spacing $\Delta t_a$ was chosen so that 20 equal intervals span the
full 30-second quench.
Within each window the residual and IC gradients remain balanced, and
accumulated error is reset at every FEM anchor.
The prediction loop is illustrated in Figure~\ref{fig:pinn_framework}.

A new network must be trained for each window.
% [HOCA #10]: "This would be problem dependent, claiming it to be <1ms might
% be misleading."
% DUZELTME: Kisitilik belirtildi: bu tezdeki ag boyutlari icin gecerli.
However, inference after training is fast.
For the network sizes used in this thesis, it takes less than
$1\,\mathrm{ms}$ on a modern graphics processing unit (GPU).
The training cost is therefore a one-time offline expense.
FEM anchor savings are realised on every subsequent simulation run.
This is precisely the property that the MSWP framework exploits.

\section{Neural Architecture Search}
\label{sec:nas_theory}

\subsection{Search Space and Strategy}

NAS automates the selection of neural network architecture \cite{zoph2018,liu2019darts}.
It defines a search space of candidate configurations and an optimisation
strategy to explore it.
Recent work has applied NAS directly to PINNs to find architectures that
minimise PDE residuals more efficiently than hand-tuned networks
\cite{wang2023,wang2025nasv2}.
In this work, the search space covers the number of hidden layers
(2 to 6), the number of neurons per layer (32 to 256), and
the activation function (ReLU, tanh, SiLU, or GELU).

\subsection{Bayesian Optimisation and TPE}

Bayesian optimisation builds a probabilistic surrogate model of the objective
function \cite{bergstra2011,snoek2012}.
It uses an acquisition function to select the next trial.
The objective here is the validation MAE.
The Tree-structured Parzen Estimator (TPE) \cite{bergstra2011}, as implemented
in Optuna \cite{akiba2019}, maintains two density models: one for good
configurations and one for bad ones.
It proposes candidates that maximise the ratio of good-to-bad likelihood.
TPE is sample-efficient and well-suited to the limited number of NAS trials
affordable in this project.

\subsection{Multi-Objective Evolutionary Search}

NSGA-II \cite{deb2002} and NSGA-III \cite{deb2014} are genetic algorithms.
They maintain a population of candidate architectures.
Each generation evolves the population using selection, crossover, and
mutation.
Both use \emph{non-dominated sorting} to identify Pareto-optimal fronts.
NSGA-II uses crowding distance to preserve diversity within a front.
NSGA-III replaces crowding distance with reference points spread uniformly
in objective space.
This makes it more effective when there are three or more objectives.
In these experiments, the two objectives are prediction MAE and model
parameter count.
The NSGA searches are implemented using \texttt{pymoo} \cite{blank2020pymoo}.

\section{The RBniCS Thermal Fin Benchmark}
\label{sec:rbnics_theory}

% [HOCA #11]: "rbniCS - ???"
% DUZELTME: RBniCS = Reduced Basis methods in computational Science.
The Reduced Basis methods in computational Science (RBniCS) library provides
a standard test problem known as the Thermal Fin benchmark \cite{rozza2008}.
The original problem models steady-state heat conduction in a two-dimensional
finned base plate with parametric conductivity.
We adapted this geometry to a three-dimensional transient quenching scenario.
The domain is a rectangular base plate with four fins extending upward.
Robin boundary conditions are applied on all surfaces.
The geometry is described in detail in Section~\ref{sec:domains}.
